\documentclass[11pt]{article}
\usepackage[a4paper,margin=24mm]{geometry}
\usepackage[T1]{fontenc}
\usepackage{lmodern,amsmath,amssymb,amsthm,mathtools,booktabs,microtype}
\usepackage[colorlinks=true,linkcolor=blue,citecolor=blue,urlcolor=blue]{hyperref}
\newcommand{\tr}{\tau}
\newcommand{\PP}{\mathsf P}
\newcommand{\sinc}{\operatorname{sinc}}
\newtheorem{theorem}{Theorem}
\newtheorem{lemma}{Lemma}
\title{Preparation memory in the complete bilocal covariance\\
\large A stationary spectral term and a neutral Gram correction}
\author{Alan Day \and Chengheng\thanks{AI research assistant (OpenAI).}}
\date{8 November 2026\quad Local research draft}
\begin{document}\maketitle
\begin{abstract}
The single-insertion projection does not preserve the full
bilocal fluctuations. Nevertheless the change from the
energy-tilted state to the actual single-selected state has a
simple two-insertion form: its uncentered part is a positive
Mehler Gram kernel built from the projected neutral tower.
All charged, neutral and mixed residual fluctuations remain
in an independently specified stationary spectral term.
The connected preparation correction includes subtraction of
the nonzero means and need not be positive. We prove its decay
under a common large center-time translation after double
scaling, and compute its leading $t^6$ and $t^7$ behavior.
\end{abstract}

\section{Scope, units and state}
Retain the overlapping complex Gaussian $p$-body CAR model with
the same original disorder, half filling, and
$p^2/N\to\lambda>0$. Exact finite-$N$ pinching precedes the
double-scaling limit. Set
\[
 C=J_s^2/\lambda,\quad q=e^{-4\lambda},\quad\zeta=\sqrt q,
 \quad h=H/\sqrt C,\quad W=\sqrt\lambda U,\quad\hbar=1.
\]
The orthogonal fields $h,U$ have unit q-Gaussian variance.
Our times are $t=\sqrt C\,t_{\rm phys}$ and
$T=\sqrt C\,T_{\rm phys}$, not the string variable
$u=2\sqrt\lambda T$. The two normalized positive states are
\begin{align}\label{states}
 \sigma(Z)&=\tr(\rho Z),&
 \rho&=1-\zeta(1+\lambda)+\zeta(h-W)^2,\\
 \tau_f(Z)&=\tr(f(h)Z),&
 f(h)&=1-\zeta+\zeta h^2.
\end{align}
The first is the actual single-selected preparation identified
in \cite{higher,sector}; the second is its energy-only restriction.
They are different states on the full neutral algebra.

Let $\mathcal A$ be the closed $h,U$ polynomial subspace in the
joint tracial $L^2$ space, and $\PP$ its orthogonal projection.
The preceding result \cite{mean} gives, for every energy-centered
single-mark descendant $X$ in the charged/neutral tower,
\begin{equation}\label{symbol}
 \PP X=\sum_{a,b}\beta_{X,ab}h^aWh^b,
 \qquad\beta_X(x,y)=\sum_{a,b}\beta_{X,ab}x^ay^b.
\end{equation}
Charged $J_d$ have $\beta_{J_d}=0$; this does not make $J_d$ zero.
Energy dressing, finite sums and the already controlled time
limits are included. Write
\[
 d\nu(x,y)=K_q(x,y)d\mu_q(x)d\mu_q(y),\quad
 (\mathcal C\beta)(x)=\int K_q(x,y)\beta(x,y)d\mu_q(y).
\]
The existing q-Mehler kernel is normalized and positive
\cite{BKS}; hence $\|\mathcal C\beta\|_2\le\|\beta\|_{L^2(\nu)}$.
The normalization \eqref{symbol} gives
$\|\PP X\|_2^2=\lambda\|\beta_X\|_{L^2(\nu)}^2$.

\newpage
\section{The two-insertion source identity}
\begin{theorem}\label{two}
For energy-centered single-mark descendants $X,Y$, including
their $h$-dressings and the controlled finite-time limits,
\begin{equation}\label{uncentered}
 \boxed{\sigma(X^\dagger Y)=\tau_f(X^\dagger Y)
 +\zeta\lambda^2(1+q)
       \langle\mathcal C\beta_X,\mathcal C\beta_Y\rangle_{\mu_q}.}
\end{equation}
The means and the connected covariance are therefore
\begin{align}\label{mu}
 \mu_X:=\sigma(X)&=-\zeta\lambda
                 \int(x+y)\beta_X(x,y)d\nu(x,y),\\
 \operatorname{Cov}_{\sigma}(X,Y)
 &=\tau_f(X^\dagger Y)+\Gamma(X,Y),\label{connected}\\
 \Gamma(X,Y)&=\zeta\lambda^2(1+q)
      \langle\mathcal C\beta_X,\mathcal C\beta_Y\rangle
                  -\overline{\mu_X}\mu_Y.\label{Gamma}
\end{align}
No replacement of $X^\dagger Y$ by $(\PP X)^\dagger\PP Y$
is made in the first term.
\end{theorem}
\begin{proof}
The state difference is
$\rho-f=\zeta[(W^2-\lambda)-\{h,W\}]$.
Use the established nilpotent source expansion for disjoint
whole-Wick $D/L$ marks \cite{multi,higher}. Each centered boundary
$W$ has no singleton source. Its only connected sources are:
another $W$ at the opposite end of the same chord, with coefficient
one, or an $L$ block containing that opposite end, also with
coefficient one. Connections to $D$ vanish. A source cannot connect
to two others, because all source cumulants of order three or
higher vanish in this class.

For two adjacent boundary $W$'s, the pairing of those two slots
is an adjacent Gaussian chord. Its contribution is exactly
$\lambda\tr(X^\dagger Y)$ and is canceled by $-\lambda$.
In every remaining source partition each boundary must attach to
a different field mark. Replacing the incident slot in each field
by $W$ gives exactly its one-$U$ projection from \cite{mean}, with
the same whole-Wick exclusions, crossing weights and chaos counts.
Thus
\begin{equation}\label{even}
 \tr[(W^2-\lambda)X^\dagger Y]
 =\tr[(W^2-\lambda)(\PP X)^\dagger\PP Y].
\end{equation}
Linear combinations and energy-regression subtractions preserve
this identity. It includes lower-chaos terms by keeping the
original field's outgoing-chord count in each diagram.

For one boundary $W$ in $\{h,W\}$, it must attach to one field.
That field is replaced by its one-$U$ term; the remaining field
can then be projected against the resulting $h,U$ test word.
Its energy regression is zero. The surviving projected word has
three $U$'s, so its trace vanishes. Therefore
\begin{equation}\label{odd}
 \tr[\{h,W\}X^\dagger Y]=0.
\end{equation}
Finally the known adjacent-four-$U$ identity \cite{response} gives
\[
 \tr[(U^2-1)f(h)Ug(h)Uk(h)]
 =(1+q)\int f(a)g(b)k(c)\,d\Pi(a,b,c),
\]
where $d\Pi=d\mu(a)K_q(a,b)d\mu(b)K_q(b,c)d\mu(c)$.
Substitution of \eqref{symbol} gives the Gram term in
\eqref{uncentered}; \eqref{mu} is the preceding one-insertion
formula. Energy centering implies $\tau_f(X)=\tau_f(Y)=0$.
\end{proof}

\newpage
\section{What remains outside the projection}
Write $X=P_X+Q_X$, $P_X=\PP X$, and similarly for $Y$.
For finite-core words, $\rho P_X^\dagger$ and
$fP_X^\dagger$ are valid $\mathcal A$ test vectors. Orthogonality
therefore kills the cross terms in either state. The formulas
extend by the bounded densities and the inherited $L^2$ limits.
Theorem~\ref{two} also applies to $Q_X,Q_Y$ by linearity,
because they are sums of single-mark terms and their projected
polynomials. Their projected symbols vanish. Consequently
\begin{equation}\label{positive}
 \operatorname{Cov}_{\sigma}(X,Y)=
 \tau_f(Q_X^\dagger Q_Y)
 +\operatorname{Cov}_{\sigma}(P_X,P_Y).
\end{equation}
Both kernels on the right are positive. The first includes the
nontrivial charged--neutral and neutral--neutral residual sectors;
their norms have not been discarded.

For example, the centered charged curvature obeys
\[
 \PP J_2=0,\qquad
 \sigma(J_2^2)=\tau_f(J_2^2)
       =4\lambda+\frac{16\lambda^2q}{1+q}>0.
\]
Substituting only projected fields would lose this entire
leading contribution to the bilocal fluctuations.

\paragraph{The connected excess is not positive in general.}
The Gram addition in \eqref{uncentered} is positive, but the
state difference in connected covariance also subtracts means.
Take the centered pure neutral size-two field
$X=S_2-\kappa_2H_2$, whose projected symbol is $\beta_X=x+y$.
Since $\mathcal C\beta_X=(1+q)x$, one obtains
\begin{align}
 \mu_X&=-2\zeta\lambda(1+q),\\
 \Gamma(X,X)&=\zeta\lambda^2(1+q)^2(1+q-4\sqrt q).
\end{align}
This is strictly negative for $q>(2-\sqrt3)^2$ on the physical
curve. It contradicts neither positivity of the full selected
variance nor \eqref{positive}. An independent rational-source
check at $q=16/81$, $\lambda=1/5$ gives exactly
\[
 \Gamma(X,X)=-\frac{1768892}{119574225},\qquad
 \operatorname{Var}_{\sigma}(X)=\frac{806826868}{773246655}>0.
\]
This rational point checks the algebra; its parameters are not
being claimed to satisfy $q=e^{-4\lambda}$.

\paragraph{A controlled scope.}
The proof concerns products of two fields, each linear in one
marked insertion before taking sums and time limits. It does not
identify a conditional expectation of arbitrary products with
the product of their projections. Four bilocal insertions have
additional internal source partitions and must retain them.

\newpage
\section{The complete bilocal two-point function}
The preceding mean calculation provides the entire symbol
\begin{align}\label{betaB}
 \beta_{T,t}(x,y)&=-2e^{iT(x-y)}\sum_{j\ge1}\Xi_j(x,y)L_j(t;x,y),\\
 L_j(t;x,y)&=\int_0^t(t-s)\sinc\!\left(\frac{(t-s)(x-y)}2\right)
              e^{-is(x+y)/2}c_j(s)\,ds.
\end{align}
Here $\PP\mathcal B(T,t)$ has W-symbol $\beta_{T,t}$,
$c_j(s)=\int e^{isx}H_j(x)d\mu_q(x)/[j]_q!$, and the
lower-chaos recursion for $\Xi_j$ is kept exactly as in \cite{mean}.
The actual double time filter and its sign are already incorporated.

For completeness, the baseline in \eqref{connected} is specified
constructively by the existing full joint spectral measure
\cite{joint,sector}, not by a new independent-probe model.
Let $X_a$ range over charged coefficients $L_n$ and neutral
coefficients $N_j$, with
\[
 k_{L_n}(s)=\frac{(is)^n}{n!},\qquad k_{N_j}(s)=2c_j(s).
\]
The common matrix measure is oriented as
\[
 d\nu_{ab}(x,y)=\tr[X_a^\dagger P_h(dx)X_bP_h(dy)],
\]
so $x$ is the output energy and $y$ the input energy. Its mixed
entries retain both $D$--$L$ connected sources and the same
disorder. Define
\[
 b_a(T,t;x,y)=-e^{iT(x-y)}\int_0^t(t-s)
 \sinc\!\left(\frac{(t-s)(x-y)}2\right)
 e^{-is(x+y)/2}k_a(s)\,ds.
\]
Then the stationary energy-tilted term is
\begin{equation}\label{base}
 \mathcal S_f(1,2)=\sum_{a,b}\int f(y)
 \overline{b_a(T_1,t_1;x,y)}b_b(T_2,t_2;x,y)\,d\nu_{ab}(x,y).
\end{equation}
It depends on the centers only through $T_2-T_1$.
All fixed-field moments of this measure are given by the existing
joint-source rules, and its absolute-continuity and time-tail
construction are inherited from \cite{joint,desc,sector}.

The result for the full actual bilocal is
\begin{equation}\label{full}
 \boxed{\operatorname{Cov}_{\sigma}(\mathcal B_1,\mathcal B_2)
 =\mathcal S_f(1,2)+\zeta\lambda^2(1+q)
 \langle\mathcal C\beta_1,\mathcal C\beta_2\rangle
                -\overline{\mu_1}\mu_2.}
\end{equation}
The density $f$ weights the \emph{input} energy in \eqref{base}.
No extra factor of $C$ appears: the physical memory's factor $C$
cancels its two physical time differentials.
The numerical results below evaluate the preparation terms and
common-center differences, not the absolute baseline integral.

\newpage
\section{Convergence and loss of preparation memory}
For the projected symbol, let $r=(1-q)^{-1/2}$,
$F_q=4\lambda q/(1-q)$, $A=qr|t|$, $B=3r|t|$. The preceding proof
gives $|c_j(s)|\le |s|^j/j!$ and
$\|\Xi_j\|_2\le4\lambda j(qr)^j+F_q(3r)^j$.
Thus a bound for both full and truncated W-symbols is
\begin{equation}
 M(t)=t^2[4\lambda A e^A+F_q(e^B-1)].
\end{equation}
For a degree cutoff $J$ satisfying $J+1>A$, $J+2>B$, set
\begin{equation}\label{tail}
 d_J(t)=t^2\left[
 \frac{4\lambda A^{J+1}}{J!\,[1-A/(J+1)]}
 +\frac{F_qB^{J+1}}{(J+1)!\,[1-B/(J+2)]}\right].
\end{equation}
Then $\|\beta-\beta_{\le J}\|_{L^2(\nu)}\le d_J(t)$,
uniformly in real $T$. Put
$c=\zeta\lambda^2(1+q)$ and
$a=\zeta\lambda\sqrt{2(1+q)}$.
The contraction of $\mathcal C$ and Cauchy--Schwarz give the
following conservative bound for the connected preparation part:
\begin{equation}\label{Gtail}
 |\Gamma(1,2)-\Gamma_{\le J}(1,2)|
 \le(c+a^2)[d_1M(t_2)+M(t_1)d_2+d_1d_2],
 \qquad d_i=d_J(t_i).
\end{equation}
These are exact-arithmetic degree estimates, excluding quadrature
and roundoff. The baseline converges by the established full-field
tails: if $\|\mathcal B_i-\mathcal B_i^{(D)}\|_2\le\delta_i$,
then its error is bounded by
\[
 \|f\|_\infty[\delta_1\|\mathcal B_2\|_2+
                         \|\mathcal B_1^{(D)}\|_2\delta_2],
 \qquad\|f\|_\infty=1-\zeta+\frac{4\zeta}{1-q}.
\]
The actual sector/preparation bridge \cite{sector} supplies the
finite-$N$ identification. Limit-only bounds are not used to
infer an unproved microscopic rate.

\begin{theorem}\label{late}
After the fixed-$\lambda$ double-scaling limit, for fixed internal
times $t_i$ and fixed center offsets $T_i$,
\begin{equation}
 \lim_{|R|\to\infty}\Gamma(T_1+R,t_1;T_2+R,t_2)=0.
\end{equation}
The selected connected two-point function consequently approaches
the stationary baseline \eqref{base} under a common translation.
\end{theorem}
\begin{proof}
A common translation multiplies $\beta_i(x,y)$ by
$e^{iR(x-y)}$. For a bounded polynomial symbol the conditional
integral in $y$ tends to zero for each $x$ by the
Riemann--Lebesgue lemma: $K_q(x,y)d\mu_q(y)$ has an integrable
density. Dominated convergence gives $L^2(\mu_q)$ convergence.
Polynomial symbols are dense in $L^2(\nu)$, and the conditional
map is a contraction for every $R$, so approximation extends the
claim to the complete symbols. The same argument applies to
$y\beta_i$ because energy support is compact. Since
$\mu_i=-\zeta\lambda\int[x\mathcal C\beta_i+
\mathcal C(y\beta_i)]d\mu_q$, the means also vanish. Equation
\eqref{Gamma} completes the proof.
\end{proof}
The order is double scaling, then large common center time.
No finite-$N$ recurrence statement, rate, or uniform joint limit
is asserted.

\newpage
\section{Short times and a full-variance difference}
The projected symbol starts as
\begin{equation}\label{shortbeta}
 \beta_{T,t}=e^{iT(x-y)}\left[
 -\frac{4i}{3}\lambda q\,t^3
 -\frac{\lambda q(1-q)}3(x+y)t^4+O_{L^2}(t^5)\right].
\end{equation}
Define the existing two-energy correlation
\[
 C_2(T)=\int e^{iT(x-y)}K_{q^2}(x,y)d\mu_q(x)d\mu_q(y),
 \qquad D_q=\frac{16}{9}\sqrt q\,\lambda^4q^2(1+q).
\]
By the Mehler semigroup,
$\|e^{iTx}P_q e^{-iT\cdot}(x)\|_2^2=C_2(T)$.
To retain the next term, use the stationary three-energy chain:
\[
 \int(x+z)e^{iT(y-z)}d\Pi(y,x,z)=\frac{i}{2}C_2'(T).
\]
Indeed its real part vanishes by simultaneous sign reversal;
subtracting its conjugate exchanges $y,z$, leaving
$\int(z-y)e^{iT(y-z)}d\Pi=iC_2'(T)$.
Substitution in \eqref{Gamma} yields
\begin{equation}\label{shortG}
 \Gamma(T,t;T,t)=D_q\left[C_2(T)t^6+
                 \frac{1-q}{4}C_2'(T)t^7\right]+O(t^8).
\end{equation}
The mean product first contributes at $t^8$. The $t^7$ term is
generally nonzero at nonzero center; the variance must not be
assumed even in the internal time in this nonstationary state.

For equal internal times the stationary baseline is independent
of common center. Therefore the following is a difference of
\emph{full} bilocal variances, not merely of projected ones:
\begin{align}\label{difference}
 &\operatorname{Var}_{\sigma}\mathcal B(T,t)
       -\operatorname{Var}_{\sigma}\mathcal B(0,t)\nonumber\\
 &\quad=\Gamma(T,t;T,t)-\Gamma(0,t;0,t)\nonumber\\
 &\quad=D_q\left[(C_2(T)-1)t^6+
                 \frac{1-q}{4}C_2'(T)t^7\right]+O(t^8).
\end{align}
The q-Hermite spectral expansion gives $0<C_2(T)<1$ for
$T\ne0$: it is a positive weighted sum of squared Fourier
coefficients, with weights $q^{2n}$, and equality at one requires
$e^{iTx}$ to be constant on the continuous energy support.
Thus the leading center-time change is negative.

The full variance itself starts at an earlier order,
\[
 \operatorname{Var}_{\sigma}\mathcal B(T,t)
 =\frac{t^4}{4}\left(4\lambda+
                   \frac{16\lambda^2q}{1+q}\right)+O(t^5),
\]
because $\mathcal B(T,t)=-t^2\alpha_T(J_2)/2+O_{L^2}(t^3)$
and the single-selected curvature restriction is stationary.
This leading term would be lost by multiplying single-field
projections. The formulas above concern ensemble/model limits
with the actual inherited preparation, not independent probes
or finite-$G$ scattering.

\newpage
\section{Executed checks and numerical scope}
\paragraph{Exact two-field checks.}
The original local CAR source evaluator with the true selected
boundary was compared against \eqref{uncentered} at
$(q,\lambda)=(1/16,1/3),(16/81,1/5)$.
It includes prefix neutral marks, charged $J_2,J_3$, the actual
ordered $N_1,\ldots,N_4$ recursion, lower-chaos/OU terms, both
field orders and inserted energy powers. The first successful
output contains 386 full-selected comparisons, 16 separate
odd-boundary zeros, 16 separate even-boundary identities,
2 charged marginal checks and 2 residual Gram matrices. All 14
principal minors of those two $3\times3$ matrices are nonnegative
as exact rationals. A separate exact witness verifies the
negative centered excess in Section~3. These formal rational
parameters test algebra and do not lie on the physical curve.

\paragraph{Fixed-$\lambda$ time diagnostics.}
On the physical curve use $q=1/4,1/2$, $\lambda=-\log q/4$,
internal times $t=1/4,1/2,1$ and centers $T=0,1/2,1,2,4$.
We compare $(J,L,n_s)=(24,48,24)$ with $(32,64,32)$, where
$L$ is the energy Jacobi dimension and $n_s$ the time quadrature
order. Scalar $c_j$ coefficients are summed starting at their
first possible Taylor order $j$; this avoids high-degree
eigensolver subtraction. Ten old mean values are reused as
read-only calibration and agree without rerunning the old driver.

Across 30 diagonal cases the maximum absolute change in
$\Gamma$ is $3.26\cdot10^{-19}$, and the maximum relative change
is $8.34\cdot10^{-14}$. The largest analytic degree bound
\eqref{Gtail} is $3.51\cdot10^{-16}$. These bounds exclude
energy/time quadrature and floating-point error; refinement
agreement is a diagnostic, not an interval certificate.
Six short-time cases at $T=1/2$ check the $t^6,t^7$ terms and
record the residual divided by $t^8$.

\begin{center}
\begin{tabular}{cccc}
\toprule
$q$&$t$&$T$&$\operatorname{Var}\mathcal B(T,t)-\operatorname{Var}\mathcal B(0,t)$\\
\midrule
$1/4$&$1/2$&$1/2$&$-3.9458574240\cdot10^{-6}$\\
$1/4$&1&$1/2$&$-2.2970480041\cdot10^{-4}$\\
$1/4$&1&1&$-5.0423075726\cdot10^{-4}$\\
$1/4$&1&2&$-6.7204451036\cdot10^{-4}$\\
\bottomrule
\end{tabular}
\end{center}
The table follows from exact cancellation of the common baseline
in \eqref{difference}; absolute total variances were not numerically
integrated here. All complex values and refinement differences
are retained in the output.

Two $6\times6$ unequal-time uncentered correction matrices are
Hermitian and numerically positive. Their connected excess
matrices instead have minimum eigenvalues about
$-8.19\cdot10^{-7}$ and $-2.21\cdot10^{-7}$, consistent with the
exact warning in Section~3. No negative eigenvalue is clipped or
reinterpreted as a failure of the full covariance's positivity.

\newpage
\section{Next physical comparison}
We now have an actual selected mean and a full two-point
reduction with the nonstationary preparation terms explicit.
The stationary baseline still carries the entire same-disorder
charged/neutral spectral matrix. This is a useful point to return
to the static-patch dictionary: compare the state, time units,
energy spectrum and candidate spatial/mass generators before
interpreting any Mehler exponent as a bulk particle mass.

Selected four-bilocal correlators require additional source
partitions and remain part of the dynamical program. The current
two-insertion identity cannot be iterated by treating products as
single-mark fields. Neither it nor its large-center-time limit
establishes finite-$G$ gravitational scattering or a complete
de Sitter holographic dual.

\begin{thebibliography}{9}\small
\bibitem{BKS} M. Bo\.{z}ejko, B. K\"ummerer and R. Speicher,
\emph{q-Gaussian processes: non-commutative and classical aspects},
Commun. Math. Phys. 185 (1997), 129--154, Sections 2 and 4,
\href{https://arxiv.org/abs/funct-an/9604010}{arXiv:funct-an/9604010}.
\bibitem{EP} E. G. Effros and M. Popa,
\emph{Feynman Diagrams and Wick Products Associated with q-Fock Spaces},
\href{https://arxiv.org/abs/math/0303045}{arXiv:math/0303045}.
\bibitem{multi} A. Day and Chengheng,
\emph{Fixed-degree multisource moments}, local proof note
\texttt{susskind-bilocal-multisource-proof}, 7 Nov. 2026,
Sections 2--6; companion \texttt{susskind-bilocal-four-point} stage.
\bibitem{higher} A. Day and Chengheng,
\emph{Higher Wick descendants in the common selected state},
local stage \texttt{susskind-slater-higher-wick}, 8 Nov. 2026.
\bibitem{joint} A. Day and Chengheng,
\emph{Charged memory: joint regression and full-time spectra},
local proof \texttt{susskind-memory-joint-regression-proof},
7 Nov. 2026, Sections 5--10.
\bibitem{desc} A. Day and Chengheng,
\emph{Ordered Wick descendants and neutral memory},
local stage \texttt{susskind-memory-wick-descendants}, 7 Nov. 2026.
\bibitem{sector} A. Day and Chengheng,
\emph{Actual fixed-charge bilocals and selected preparations},
local stage \texttt{susskind-sector-bilocal}, 8 Nov. 2026.
\bibitem{response} A. Day and Chengheng,
\emph{Preparation memory in the neutral slope}, local stage
\texttt{susskind-neutral-selected-response}, 8 Nov. 2026.
\bibitem{mean} A. Day and Chengheng,
\emph{A two-energy formula for the selected bilocal mean},
local stage \texttt{susskind-selected-bilocal-mean}, 8 Nov. 2026.
\end{thebibliography}
\end{document}
